\clearpage\SourceLeaf{369}\ArithmeticHead{362}{L}{ANALYSIS.\quad BRIEFWECHSEL.}
\NumberedItem{3}
\HistoricTitle{\textsc{Gauss} an \textsc{Bessel}. Göttingen, 21. November 1811.}
\HistoricCitation{[Briefwechsel zwischen \textsc{Gauss} und \textsc{Bessel}, Leipzig 1880, S. 151--154.]}
\noindent {.}\quad{.}\quad{.}\quad{.}\quad{.}\quad{.}\quad{.}\quad Mit Vergnügen sehe ich aus Ihrem Briefe, dass Sie auch über die Facultäten gearbeitet haben. Wahrscheinlich werden sich unsre Ansichten öfters begegnen. Ich arbeite nemlich jetzt an einer Abhandlung für unsere Soc[ietät], die in etwa 6 Wochen vollendet seyn wird, und die Reihe
{\small
\[
1+\frac{\alpha{.}\beta}{1{.}\gamma}x+\frac{\alpha{.}\alpha+1{.}\beta{.}\beta+1}{1{.}\,2{.}\gamma{.}\gamma+1}xx+\frac{\alpha{.}\alpha+1{.}\alpha+2{.}\beta{.}\beta+1{.}\beta+2}{1{.}\,2{.}\,3{.}\gamma{.}\gamma+1{.}\gamma+2}x^3+\text{etc.}
\]
}
\noindent betrifft, und auch die Functionen, die mit \textsc{Kramps} Facultäten[\textsuperscript{*})] zusammenhangen, berührt. Eine solche Function ist allerdings in der Analyse einzuführen; ich finde es aber nicht zweckmässig, eine solche, wie \textsc{Kramp} gethan hat, $a^{bIc}$, die von drei veränderlichen Grössen abhängt aufzunehmen, da man mit einer Function Einer veränderlichen Grösse $1^{xI1}$ vollkommen eben so weit reicht. Also das Product
\[
1{.}\,2{.}\,3\,{.}\,{.}\,{.}\,x=\Pi x
\]
\noindent ist die Function, die meiner Meinung nach in der Analyse eingeführt werden muss, und auf die so wie auf die davon derivirten Functionen
\[
\frac{d\Pi x}{\Pi x{.}\,dx}\ \text{etc.}
\]
\noindent unzählige Aufgaben sich zurückführen lassen. Will man sich aber nicht wie \textsc{Kramp} zahllosen Paralogismen und Paradoxen und Widersprüchen[\textsuperscript{**})] blossstellen, so muss $1{.}\,2{.}\,3\,{.}\,{.}\,{.}\,x$ nicht als \textls[120]{Definition} von $\Pi x$ gebraucht werden, da eine solche nur, wenn $x$ eine ganze Zahl ist, einen bestimmten Sinn hat, sondern man muss von einer höheren \textls[120]{allgemein}, selbst auf imaginäre Werthe von $x$ anwendbaren, Definition ausgehen, wovon, wie Sie mit Recht bemerken,
\par\medskip\noindent\rule{.24\linewidth}{.4pt}\par\smallskip
{\small
\noindent\hspace*{\parindent}[\textsuperscript{*}) Siehe: \textit{Analyse des réfractions astronomiques et terrestres}, par le citoyen \textsc{Kramp}, Leipzig 1799, Chapitre III, Analyse des facultés numériques, S. 46 ff.]
\par\noindent\hspace*{\parindent}[\textsuperscript{**}) Siehe z. B. die von \textsc{Bessel}, \textit{Briefwechsel}, S. 200, erwähnten Stellen in \textsc{Kramps} Abhandlung: § 40, S. 56 und § 200, S. 108.]
\par}
\clearpage\SourceLeaf{370}\ArithmeticHead{363}{R}{GAUSS AN BESSEL.}
\noindent jene als specieller Fall erscheint. Ich habe folgende gewählt[\textsuperscript{*})]:
\[
\Pi x=\frac{1{.}\,2{.}\,3\,{.}\,{.}\,{.}\,k{.}\,k^x}{x+1{.}\,x+2{.}\,x+3\,{.}\,{.}\,{.}\,x+k},
\]
\noindent wenn $k$ unendlich wird. Ich freue mich, dass unsre Ansichten hier wie in andern Puncten zusammentreffen.

\indent Das \textls[120]{Wesen} einer Function ist das \textls[120]{Bestimmtseyn}, und darauf gründet sich das Haupterforderniss einer guten Definition. Höchst verwerflich ist daher eine Definition aus einer ihrer Natur nach divergierenden Reihe wie \textsc{Kramp}’s $\Gamma x$[\textsuperscript{**})], welches
\[
=\log\Pi\frac{1}{x}+\left(\frac{1}{x}+\frac{1}{2}\right)\log x+\frac{1}{x}-\frac{1}{2}\log2\pi.
\]
\noindent Eine Reihe die nicht immer convergirt, wie meine obige, kann auch nur innerhalb der Schranken, wo sie convergirt, als Definition gelten, und die Frage was aus der Function werden soll\textsuperscript{***}), wenn man ihr Argument jene Schranken überschreiten lässt, verträgt nur dann erst eine Antwort, wenn ein höheres
\par\medskip\noindent\rule{.24\linewidth}{.4pt}\par\smallskip
{\small
\noindent\hspace*{\parindent}[\textsuperscript{*}) In Wirklichkeit geht diese Definition der Funktion $\Pi x$ auf \textsc{Euler} zurück, der sich in dem Briefe an \textsc{Goldbach} vom 13. Oktober 1729 (P. H. \textsc{Fuss}, \textit{Correspondance mathématique et physique} I, St Pétersbourg 1843, S. 3) und in den \textit{Institutiones calculi differentialis} II, 1755, cap. XVII, Opera omnia ser. I, vol. 10, S. 641 zur Definition von $x!$ des unendlichen Produkts
\[
\frac{1^{1-x}{.}\,2^x}{1+x}\cdot\frac{2^{1-x}{.}\,3^x}{2+x}\cdot\frac{3^{1-x}{.}\,4^x}{3+x}\ \mathord{\cdot}\,\mathord{\cdot}\,\mathord{\cdot}
\]
\noindent bedient. \textsc{Gauss} weist (\textit{Disquisitiones circa seriem}, art. 20, Werke III, S. 145, 146) selbst — allerdings ohne \textsc{Euler} zu nennen — auf die Identität seiner Definition von $\Pi x$ mit der von \textsc{Euler} gegebenen hin.]
\par\noindent\hspace*{\parindent}[\textsuperscript{**}) Siehe a. a. O. § 179, S. 102, wo $\Gamma y$ durch die Reihe
\[
\beta y+\tfrac{1}{3}\delta y^3+\tfrac{1}{5}\zeta y^5+\tfrac{1}{7}\vartheta y^7+\mathord{\cdot}\,\mathord{\cdot}\,\mathord{\cdot}
\]
\noindent erklärt wird, in der nach § 129, 130, S. 85
\[
\beta=\tfrac{1}{1\,2},\quad\delta=-\tfrac{1}{1\,2\,0},\quad\zeta=\tfrac{1}{2\,5\,2},\quad\vartheta=-\tfrac{1}{2\,4\,0},\,{.}\,{.}\,{.}
\]
\noindent also
\[
\beta=\tfrac{1}{2}\mathfrak A,\quad\delta=-\tfrac{1}{4}\mathfrak B,\quad\zeta=\tfrac{1}{6}\mathfrak C,\quad\vartheta=-\tfrac{1}{8}\mathfrak D,\,{.}\,{.}\,{.},
\]
\noindent wo $\mathfrak A,\mathfrak B,\mathfrak C,\mathfrak D,\,{.}\,{.}\,{.}$ die \textsc{Bernoulli}schen Zahlen sind; vergl. den art. 20. der \textit{Disquisitiones circa seriem}, Werke III, S. 152.]
\par}
\noindent\hspace*{\parindent}\textsuperscript{***}) Man sollte überhaupt nie vergessen, dass die Functionen, wie alle mathematischen Begriffszusammensetzungen, nur unsere eignen Geschöpfe sind, und dass, wo die Definition, von der man ausging, aufhört einen Sinn zu haben, man eigentlich nicht fragen soll \textls[120]{was ist?} sondern was \textls[120]{convenirt anzunehmen?} damit ich immer consequent bleiben kann. So z. B. das Product aus $-\mathord{\cdot}-.$
\par\smallskip\hfill46\textsuperscript{*}\par
\clearpage\SourceLeaf{371}\ArithmeticHead{364}{L}{ANALYSIS.\quad BRIEFWECHSEL.}
\noindent Princip zur Definition aufgefunden ist, wie bei meiner Reihe die Differentialgleichung
\[
0=\alpha\beta P-(\gamma-(\alpha+\beta+1)x)\frac{dP}{dx}-(x-xx)\frac{ddP}{dx^2}\ \text{*)}.
\]
\noindent Auf gleiche Art fallen alle Paradoxa, die einige Mathematiker bei den Logarithmen gefunden haben, von selbst weg, wenn man nicht von der gewöhnlichen Definition $\mathrm{basis}^{\mathrm{logar.}}=\text{Zahl}$ ausgeht, die eigentlich nur dann recht befriedigend ist, wenn der Exponent eine ganze Zahl ist, und ganz [und gar] keinen Sinn gibt, wenn der Exponent gar imaginär werden soll — sondern jede Grösse die in der Reihe
\[
1+x+\tfrac{1}{2}xx+\tfrac{1}{6}x^3+\text{etc.}\,{.}\,{.}\,{.}
\]
\noindent für $x$ subst[ituirt] ihr den Werth $A$ gibt, Logarithm von $A$[\textsuperscript{**})] \textls[120]{nennt}; wie ich mit Vergnügen sehe, hat der Portugiese \textsc{Acunha}[\textsuperscript{***})] diese Definition wirklich gewählt — in einer schlechten Rec[ension] in unseren Gel[ehrten] Anz[eigen $\dagger$)] ist er deswegen getadelt — aber freilich die nöthigen Beweise, dass allgemein
\[
(1+x+\tfrac{1}{2}xx\,{.}\,{.}\,{.})(1+y+\tfrac{1}{2}yy\,{.}\,{.}\,{.})=1+(x+y)+\tfrac{1}{2}(x+y)^2+{.}\,{.}\,{.}
\]
\noindent nur schlecht, nemlich bloss durch Induction ausgeführt, obgleich es \textls[120]{so sehr} leicht streng zu beweisen ist.

\indent So wie es mir immer Vergnügen macht, in Wesentlichen Dingen mit Ihnen, lieber \textsc{Bessel}, zusammen zu treffen, so kann ich nicht leugnen, dass ich auch mit Vergnügen bemerke, wenn wir in Kleinigkeiten, die mehr nur Sache des Geschmacks sind, uns begegnen. Noch heute las ich ihre Recension
\par\medskip\noindent\rule{.24\linewidth}{.4pt}\par\smallskip
\noindent\hspace*{\parindent}\textsuperscript{*}) Nach welcher \textls[120]{jedem}, reellen oder imaginären Werthe von $x$ vermittelst eines \textls[120]{stetigen} Überganges von $x=0$, ohne $x=1$ zu berühren (was durch imaginäre Zwischenwerthe geschehen kann), ein, gewöhnlich mehrere, ja unendlich viele Werthe von $P$ entsprechen.
\par{\small
\noindent\hspace*{\parindent}[\textsuperscript{**}) In der Handschrift steht $x$ statt $A$.]
\par\noindent\hspace*{\parindent}[\textsuperscript{***}) \textsc{José Anastasio da Cunha}, \textit{Principios Mathematicos}, Lisboa 1790. \textsc{Gauss} kannte wahrscheinlich die Übersetzung \textit{Principes Mathématiques} de feu \textsc{J.-A. da Cunha}, traduits littéralement du Portugais par \textsc{J. M. d’Abreu}, Bordeaux 1811, siehe Livre IX, S. 117 ff., wo in dem von \textsc{Gauss} angegebenen Sinne von den Potenzen und Logarithmen gehandelt wird.]
\par\noindent\hspace*{\parindent}[$\dagger$) Göttingische Gelehrte Anzeigen 1811, 181. Stück, S. 1801, siehe besonders S. 1804--1805; diese Anzeige bezieht sich auf die französische Übersetzung des in Rede stehenden Werks.]
\par}
\clearpage\SourceLeaf{372}\ArithmeticHead{365}{R}{GAUSS AN BESSEL.}
\noindent von \textsc{Cagnoli} in der Jen[aischen] A[llgemeinen] L[iteraturzeitung\textsuperscript{*})], wo es mich freute, dass Sie sich gegen das widerwärtige $\sim$ erklären. Auch ist mir jedesmal fatal das garnicht analogische $\sin^2\varphi$, obgleich auch \textsc{Laplace} es gebraucht; fürchtet man, dass $\sin\varphi^2$ zweideutig werden könne (was doch vielleicht nie oder höchst selten eintritt, wenn man von $\sin(\varphi^2)$ spräche), ei nun, so schreibe man $(\sin\varphi)^2$ aber nicht $\sin^2\varphi$, was der Analogie nach nur $\sin(\sin\varphi)$ bedeuten sollte. Noch unausstehlicher ist mir das barbarische den Franzosen nachgeäffte »\textls[120]{Sey} $a=b$«. Es sind diess freilich alles nur nichtswürdige Kleinigkeiten, aber wenn mir bei einer analytischen Ausführung Eleganz immer ein wohlthuender Genuss ist, so mag ich auch gern, dass am äusseren Kleide nichts mein Auge beleidigt.\quad {.}\quad{.}\quad{.}\quad{.}\quad{.}\quad{.}\quad{.}\quad{.}\quad{.}\quad{.}\quad{.}\quad{.}\quad{.}\quad{.}\quad{.}\quad{.}\quad{.}\quad{.}\quad{.}
\par\medskip\noindent\rule{.24\linewidth}{.4pt}\par\smallskip
{\small
\noindent\hspace*{\parindent}[\textsuperscript{*}) Jahrgang 1811, Nr. 206, 207, \textit{Recensionen von Fr. W. Bessel}, herausgegeben von \textsc{Engelmann}, Leipzig 1878, S. 138; die Rezension bezieht sich auf die \textit{Trigonométrie rectiligne et sphérique} par \textsc{Antoine Cagnoli}, traduite de l’Italien par \textsc{N. M. Chompré}, 2\textsuperscript{de} éd., Paris 1808.]
\par}
